\section{Conclusions}
\label{sec:conclusions}

This project compared regular global \texttt{new} with overloaded global
\texttt{new} backed by eight versions of CustomMemoryAllocator.
We ran four workloads at three time scales across several machines, and we
measured elapsed time, peak memory, page faults, and context switches.

\subsection{No custom version was faster everywhere}

Version v7 lost to the system allocator in all twelve combinations of workload
and scale on the Ryzen host.
It was about 9.7 percent slower on the matrix array, 9.7 percent slower on
uniform nodes, 4.5 percent slower on nested stress, and 12 percent slower on
the linked list.
Version v6 on the Apple M3~Max did better, but only on two workloads.
It improved the matrix array by 2.3 percent and uniform nodes by 3.1 percent,
while losing 6.4 percent on nested stress and matching the system allocator on
the linked list.
In other words, the workload decided the outcome more than the allocator
version did.

\subsection{Newer versions were not better versions}

The matrix campaign ran all seven released versions on one machine with one set
of parameters.
The two simplest allocators won.
Version v1 used a single address-ordered free list with first-fit search and
finished 3.8 percent faster than the system allocator.
Version v2 was nearly identical.
From there the results got worse rather than better.
Version v5 was 2.7 percent slower than regular \texttt{new} and version v7 was
1.8 percent slower.
Each of those versions added machinery that was meant to help, such as
deferred coalescing and adaptive per-thread caches.
On this workload that machinery only added cost.

\subsection{The custom allocator traded memory for system calls}

The pool asks the operating system for memory once and never gives it back.
That choice showed up clearly in the measurements.
The overloaded programs took 97 to 99 percent fewer minor page faults than
regular \texttt{new} and spent roughly ten times less time in the kernel.
However, the savings did not cover the cost of the allocator's own
bookkeeping.
Peak memory tells the same story from the other side.
The overhead was about 4 percent on the matrix array and 3 percent on uniform
nodes, but it reached 28 percent on nested stress and 155 percent on the
linked list.

\subsection{Our improved allocator succeeded on its target and failed
elsewhere}

Version v8 started from v7 and changed the paths that matter for packed matrix
rows.
It beat regular \texttt{new} by 4.5 to 5.1 percent on the matrix array, and the
result was statistically significant at every scale.
The same changes made the linked list unusable.
Those runs had to be killed before finishing.
This is the clearest evidence in the project for the point that runs through
all of our results.
An allocator is tuned for an assumed pattern of requests, and it pays for that
assumption whenever the pattern does not hold.

\subsection{The environment mattered too}

The same v7 code was about 5 percent slower than regular \texttt{new} on
Group~1's Core i5-6400 and about 11 percent slower on our Ryzen host.
The pthreads experiments showed the same sensitivity
(\Cref{sec:pthreads-experiments}).
The custom allocator hurt performance under Windows WSL but was statistically
indistinguishable from regular \texttt{new} under native Linux.
Results from one machine should not be treated as a general claim.

\subsection{What we conclude}

A general-purpose system allocator is hard to beat because it has been tuned
for mixed workloads over decades.
A custom allocator is worth building when profiling shows that allocation is on
the critical path and the request pattern is known in advance.
Our measurements support that conclusion in both directions.
The custom versions won on the workload whose allocation pattern was regular
and predictable, and they lost badly on the workload full of small,
short-lived, irregular objects.
